% Flavor catalog per-process template.
% process_id: CR009
% family: collider_rs
% owner_agent_id: PKA-CR009
% writer_agent_id: null
% checker_agent_id: null
% opus_signoff_id: null
% Canonical metadata sidecar: CR009.yaml
% Status is controlled by the YAML sidecar status_history, not this comment.

\section{\texorpdfstring{\(pp\to \ell^+\ell^-\) High-Mass Tail Contact-Operator Bound}{pp to l+l- High-Mass Tail Contact-Operator Bound}}

\subsection*{Process}
This entry covers the neutral-current Drell--Yan high-mass tail,
\(pp\to \ell^+\ell^-+X\) with \(\ell=e,\mu\), interpreted as a
model-independent four-fermion contact interaction between light quarks and
charged leptons.  The collider signature is a smooth excess or deficit in the
dilepton invariant-mass distribution rather than a narrow resonance.

\subsection*{Current best limit(s)}
The current PDG summary of quark and lepton compositeness quotes the
full-Run-2 ATLAS dielectron+dimuon contact-interaction limits as the highest
single-experiment reach in the conventional \(g_{\rm contact}^2=4\pi\)
normalization.  For the left-left model, ATLAS obtains
\(\Lambda_{\rm LL}^{+}>35.8~\mathrm{TeV}\) for constructive interference and
\(\Lambda_{\rm LL}^{-}>26.0~\mathrm{TeV}\) for destructive interference at
95\% CL.  The same PDG paragraph gives \(\Lambda_{\rm RR}^{+}>35.5~\mathrm{TeV}\),
\(\Lambda_{\rm RR}^{-}>26.5~\mathrm{TeV}\), \(\Lambda_{\rm LR}^{+}>32.5~\mathrm{TeV}\),
and \(\Lambda_{\rm LR}^{-}>28.8~\mathrm{TeV}\)
(\texttt{pdg2025\_compositeness\_extract.txt}).

CMS provides the independent full-Run-2 cross-check using up to
140~\(\mathrm{fb}^{-1}\) at 13 TeV.  Its published contact-interaction
limits range from about \(23.8\)--\(23.9~\mathrm{TeV}\) to
\(36.4~\mathrm{TeV}\), depending on the helicity current and interference
sign (\texttt{cms\_2021\_arxiv2103\_02708.txt}).  The HEPData record supplies
the structured limit tables and the binned inputs needed for any serious
reinterpretation (\texttt{cms\_hepdata\_2021\_ci\_limits\_extract.txt}).

\subsection*{Relevance to the RS / anarchic-flavor pipeline}
In a 5D RS interpretation, this is primarily a constraint on neutral heavy
vector exchange that induces semileptonic operators of the form
\((\bar q\gamma_\mu q)(\bar\ell\gamma^\mu\ell)\).  It is not a direct
KK-gluon mass limit, since a color-octet KK gluon does not mediate
\(\ell^+\ell^-\) at tree level.  It constrains the electroweak KK tower, a
\(Z^\prime\)-like custodial vector, or a more generic compositeness sector
only after specifying the light-quark and lepton couplings.

With the usual contact-interaction convention,
\(4\pi/\Lambda^2 \sim |g_q g_\ell|/M_V^2\).  Therefore the ATLAS
\(\Lambda_{\rm LL}^{+}>35.8~\mathrm{TeV}\) limit corresponds only
schematically to
\[
  M_V \gtrsim \sqrt{|g_q g_\ell|/(4\pi)}\,35.8~\mathrm{TeV},
\]
not to a universal \(35.8~\mathrm{TeV}\) resonance mass bound.  For
\(|g_q g_\ell|\sim 1\), this is an \(O(10~\mathrm{TeV})\) neutral-vector
scale; for strongly coupled light fermions it can be higher, while for
UV-localized light fermions it is weaker.  The quark-scan methodology note
already finds the low-energy anarchic-flavor median scale
\(\Mkk^{\min,p50}=47.26~\mathrm{TeV}\) at \(\gs=3\), with a
95\%-acceptance crossing at \(127.13~\mathrm{TeV}\).  CR009 is therefore a
cross-check and a handle on non-anarchic or non-universal RS variants, not the
leading anarchic-flavor constraint.

\subsection*{Post-2010 developments}
The LHC history is a steady increase in energy, luminosity, and analysis
sophistication.  ATLAS started with 7 TeV dielectron and dimuon data in
arXiv:1112.4462, reaching \(O(10~\mathrm{TeV})\) in the LL isoscalar
benchmark.  ATLAS then used the 8 TeV dilepton spectrum and
forward-backward asymmetry in arXiv:1407.2410, raising the contact-scale
range to \(15.4\)--\(26.3~\mathrm{TeV}\).  CMS's 8 TeV spectrum analysis
in arXiv:1412.6302 provided comparable same-flavor contact limits and
model-independent resonance information.

Run 2 changed the reach.  ATLAS's first 13 TeV result with 3.2
\(\mathrm{fb}^{-1}\), arXiv:1607.03669, already reached
\(16.7\)--\(25.2~\mathrm{TeV}\).  The 36.1 \(\mathrm{fb}^{-1}\) ATLAS
search, arXiv:1707.02424, pushed the model-dependent range to
\(24\)--\(40~\mathrm{TeV}\).  CMS's 36 \(\mathrm{fb}^{-1}\) result,
arXiv:1812.10443, reached \(20\)--\(32~\mathrm{TeV}\).  The current
full-Run-2 anchors are ATLAS arXiv:2006.12946 and CMS arXiv:2103.02708.

The SMEFT literature clarified how to use these tails.  Falkowski,
Gonzalez-Alonso, and Mimouni compiled low-energy four-fermion constraints in
arXiv:1706.03783, while later high-energy Drell--Yan SMEFT work treats the
tails as energy-growing probes of semileptonic dimension-six operators
(\texttt{smeft\_theory\_arxiv\_extracts.txt}).

\subsection*{Constraint validity and model dependence}
The quoted \(\Lambda\) limits are measured exclusions, but the map to RS is
not model-independent.  The experimental interpretation assumes a particular
helicity current, an interference sign, and the conventional
\(g_{\rm contact}^2=4\pi\) normalization.  The CMS combination assumes a
universal contact interaction for electrons and muons.  ATLAS and CMS also
use different background strategies, and the high-mass tail is sensitive to
PDF and electroweak corrections.

For RS, the largest ambiguity is the light-fermion coupling pattern.  A
minimal light-fermion UV localization suppresses the contact operator, while
non-anarchic lepton localization, partial compositeness of light quarks, or an
extra neutral vector with sizable dilepton branching can make the bound much
more relevant.  EFT validity also requires the mediator to be heavier than the
dilepton masses driving the fit; otherwise the correct object is a resonance
or broad-resonance likelihood, not the contact-limit number alone.

\subsection*{Code coverage in this repo}
\textbf{NO}.  The required grep over \texttt{quarkConstraints/}, \texttt{qcd/},
\texttt{flavorConstraints/}, \texttt{neutrinos/}, \texttt{yukawa/},
\texttt{warpConfig/}, \texttt{solvers/}, \texttt{scanParams/}, and
\texttt{tests/} found no Drell--Yan, dilepton-tail, contact-interaction,
collider-reinterpretation, CheckMATE, MadAnalysis5, SModelS, Delphes,
Pythia, or MadGraph implementation.  The only \texttt{CMS} match was the
unrelated RunDec alpha-s test at \texttt{tests/test\_alpha\_s.py:88}.

Adjacent evidence shows that the scan computes \(M_{KK}\) and low-energy
\(\Delta F=2\) ratios, not collider direct-search filters:
\texttt{quarkConstraints/scan.py:359} maps \(\Lambda_{\rm IR}\) to \(M_{KK}\),
\texttt{quarkConstraints/scan.py:376--380} evaluates flavor diagnostics and
KK-gluon \(\Delta F=2\) constraints, and
\texttt{quarkConstraints/deltaf2.py:320--324} states that the implemented
Wilsons are tree-level KK-gluon-inspired \(\Delta F=2\) coefficients matched
at \(M_{KK}\).

\subsection*{Implementation difficulty}
\textbf{HIGH}.  Recording the bound is easy; using it as a live scan
constraint is not.  A faithful implementation would need either a public
likelihood or a recreated binned dilepton-tail likelihood with correlations,
PDF and electroweak systematics, an RS-to-SMEFT matching layer for the
neutral electroweak KK vectors, and validity logic deciding when a point
should be treated as a contact interaction rather than a resolved resonance.
External reinterpretation infrastructure such as MadAnalysis5, CheckMATE, or
a dedicated high-\(p_T\) SMEFT tool would be the pragmatic route.

\subsection*{Key references}
\texttt{PDG2025\_QuarkLeptonCompositeness};
\texttt{ATLAS2020\_NonResonantDilepton};
\texttt{CMS2021\_HighMassDilepton};
\texttt{CMS2021\_HEPData\_CILimits};
\texttt{FalkowskiGonzalezAlonsoMimouni2017\_SMEFT4F};
\texttt{DawsonGiardinoIsmail2018\_DYSMEFT}.
